បណ្ឌិត្យសភាបច្ចេកវិទ្យាឌីជីថលកម្ពុជា (CADT) សូមជម្រាបជូនដំណឹងដល់សិស្សានុសិស្សឱ្យបានជ្រាបថា ក្រោមការយកចិត្តទុកដាក់ជាប្រចាំលើអនាគតយុវជន យុវនារី ក្នុងយុគសម័យឌីជីថល តាមរយៈក្រសួងប្រៃសណីយ៍ និងទូរគមនាគមន៍ (ក.ប.ទ) អាហារូបករណ៍ទេពកោសល្យឌីជីថលតេជោ ត្រូវបានបង្កើតឡើងចាប់តាំងពីឆ្នាំ២០២១ មក។ ក្រោមការដឹកនាំប្រកបដោយគតិបណ្ឌិតរបស់ សម្តេចមហាបវរធិបតី ហ៊ុន ម៉ាណែត នាយករដ្ឋមន្ត្រី នៃព្រះរាជាណាចក្រកម្ពុជា ក្នុងឆ្នាំសិក្សាថ្មី ២០២៥-២០២៦នេះ កម្មវិធីអាហារូបករណ៍ទេពកោសល្យឌីជីថលតេជោ នឹងបន្តផ្តល់អាហារូបករណ៍សម្រាប់សិក្សាថ្នាក់បរិញ្ញាបត្រជំនាញបច្ចេកវិទ្យាឌីជីថលនៅបណ្ឌិត្យសភា CADT ចំនួន ២០០ កន្លែង ក្នុងគោលបំណងផ្តល់ឱកាសឱ្យនិស្សិតកាន់តែច្រើនថែមទៀតអាចសិក្សាជំនាញបច្ចេកវិទ្យាឌីជីថល ដែល ជាជំនាញមួយមានតម្រូវការខ្ពស់ក្នុងទីផ្សារការងារបច្ចុប្បន្ន ហើយចាំបាច់ក្នុងការអភិវឌ្ឍសេដ្ឋកិច្ច និងសង្គមឌីជីថល នៅកម្ពុជា។\\%
\CADTNumberItem{១. អត្ថប្រយោជន៍នៃអាហារូបករណ៍}
\CADTBodyParagraph{និស្សិតជ័យលាភីអាហារូបករណ៍ទេពកោសល្យឌីជីថលតេជោនឹងទទួលបានអត្ថប្រយោជន៍រួមមាន៖}
\CADTDashItem{កុំព្យូទ័រយួរដៃថ្មីចំនួន ១គ្រឿង}
\CADTDashItem{ទទួលបានការរ៉ាប់រងថ្លៃសិក្សា១០០\% រហូតដល់ចប់ថ្នាក់បរិញ្ញាបត្រដែលមានតម្លៃស្មើ ៤០.០០០.០០០៛ (សែសិបលានរៀល)}
\CADTDashItem{ទទួលបានឱកាសការងារល្អ និងប្រាក់បៀវត្សខ្ពស់ក្នុងក្រសួង ស្ថាប័ន ក្រុមហ៊ុន ឬធនាគារជាដើម។}
\CADTNumberItem{២. ជំនាញសិក្សា}
\CADTBodyParagraph{អាហារូបករណ៍ទេពកោសល្យឌីជីថលតេជោផ្តល់ជូនដល់និស្សិតសម្រាប់សិក្សានៅបណ្ឌិត្យសភា CADT ដែលមានជំនាញ និងចំនួនកំណត់ដូចខាងក្រោម៖}
\CADTDashItem{បរិញ្ញាបត្រជំនាញទូរគមនាគមន៍ និងបណ្តាញ (Bachelor of Telecommunications \& Networking)}
\CADTPlusSeatItem{ឯកទេស សន្តិសុខសាយប័រ (Specialized in Cyber Security)}{ចំនួន ៤០កន្លែង}
\CADTPlusSeatItem{ឯកទេស វិស្វកម្មទូរគមនាគមន៍ និងបណ្តាញ (Specialized in Telecommunications \& Network Engineering)}{ចំនួន ៤៥កន្លែង}
\CADTDashItem{បរិញ្ញាបត្រជំនាញវិទ្យាសាស្ត្រកុំព្យូទ័រ (Bachelor of Computer Science)}
\CADTPlusSeatItem{ឯកទេស វិស្វកម្មសុសវែរ (Specialized in Software Engineering)}{ចំនួន ៤០កន្លែង}
\CADTPlusSeatItem{ឯកទេស វិទ្យាសាស្ត្រទិន្នន័យ (Specialized in Data Science)}{ចំនួន ៣៥កន្លែង}
\CADTDashItem{បរិញ្ញាបត្រជំនាញធុរកិច្ចឌីជីថល (Bachelor of Digital Business)}
\CADTPlusSeatItem{ឯកទេស ពាណិជ្ជកម្មអេឡិចត្រូនិក (Specialized in e-Commerce)}{ចំនួន ៤០កន្លែង}
\CADTNumberItem{៣. លក្ខខណ្ឌនៃការជ្រើសរើស}
\CADTBodyParagraph{បេក្ខជន-បេក្ខនារីទាំងអស់ត្រូវឆ្លងកាត់ការប្រឡងប្រជែងដែលរៀបចំដោយបណ្ឌិត្យសភា CADT ដោយមានលក្ខខណ្ឌដូចខាងក្រោម៖}
\CADTDashItem{ជាសិស្សជាប់សញ្ញាបត្រមធ្យមសិក្សាទុតិយភូមិឆ្នាំសិក្សា២០២៥ មាននិទ្ទេសរួម A B ឬ C}
\CADTDashItem{ឆ្លងកាត់ការប្រឡងប្រជែងលើវិញ្ញាសា គណិតវិទ្យា តក្កវិទ្យា និងភាសាអង់គ្លេស}
\CADTDashItem{ការប្រឡងធ្វើឡើងចំនួន ២ដង តាមរយៈប្រព័ន្ធកុំព្យូទ័រ និងធ្វើការសម្ភាសន៍ដើម្បីពិនិត្យលើឆន្ទៈ។ សិស្សដែលមានជីវភាពក្រីក្រមកពីតំបន់ដាច់ស្រយាល ជនមានពិការភាព និងសិស្សនារីត្រូវបាន លើកទឹកចិត្ត។}
\CADTNumberItem{៤. កាលបរិច្ឆេទ និងកន្លែងទទួលពាក្យ}
\CADTBodyParagraph{ការចុះឈ្មោះប្រឡងតាមប្រព័ន្ធអនឡាញ៖}
\CADTDashItem{និស្សិតចុះឈ្មោះតាមរយៈតំណភ្ជាប់ \url{https://cadt.edu.kh/techo} ចាប់ពីថ្ងៃជូនដំណឹងនេះ រហូតដល់ ថ្ងៃទី២២ ខែតុលា ឆ្នាំ២០២៥}
\CADTDashItem{បន្ទាប់ពីបេក្ខជន-បេក្ខនារី ដែលបានចុះឈ្មោះប្រឡងតាមប្រព័ន្ធ online រួច នឹងត្រូវមកទទួល សលាកបត្រសុំសិទ្ធិប្រឡងនៅអគារបណ្ឌិត្យសភាបច្ចេកវិទ្យាឌីជីថលកម្ពុជា}
\CADTDashItem{ការប្រឡងនឹងធ្វើឡើ...អគារបណ្ឌិត្យសភាបច្ចេកវិទ្យាឌីជីថលកម្ពុជារៀងរាល់ម៉ោងធ្វើការចាប់ពី ថ្ងៃចន្ទ ដល់ ថ្ងៃសុក្រ។}
\CADTNumberItem{៥. ទំនាក់ទំនងព័ត៌មានតាមរយៈ}
